\section{Research Objectives and Questions}
\label{sec:questions}

We investigate a generator that maps a set of prompt summaries to one shared
adapter for a frozen EAGLE3 drafter. Separate serving studies compare the
original drafter, trained static LoRA,
shared-adapter generators and per-request generators. Checkpoint labels and
protocol changes are retained within each study rather than pooled. Historical training
ablations are recorded separately: fixed-update data scaling and the decision
to discontinue bank-based generator warmup.

The working questions are:
\begin{enumerate}[leftmargin=*]
  \item How sensitive is generator validation behavior to training-pool size
  under a fixed optimizer-update budget?
  \item How do generator checkpoints behave relative to the original drafter and
  a static trained adapter under a specified serving workload?
  \item How do recorded measurements differ between Nemotron validation prompts
  and selected Red Hat benchmark categories?
  \item What changes when one generated adapter is shared by a fixed batch of
  requests at different concurrencies?
  \item Does conditioning on more examples improve a shared adapter when
  training data and optimizer updates are matched?
  \item Do per-request adapters preserve acceptance while permitting native
  continuous admission, and does better acceptance yield higher throughput?
  \item How does LoRA rank relate to adapter variation, category separation,
  acceptance and serving token rate?
  \item What time and memory costs arise from obtaining and reusing target
  activations during training?
\end{enumerate}

These questions remain open. A numerical table records an observation under its
measurement conditions; it does not establish general superiority, statistical
significance, lossless acceleration, or completion of a planned experiment.

\subsection{Record Conventions}
Experiment identifiers remain stable as evidence is appended. We use
\emph{planned}, \emph{implemented}, \emph{partially measured}, \emph{measured}, and
\emph{investigation deferred} as execution statuses. Evidence quality is stated
separately: a measured value may still have incomplete provenance or no repetition.
Dates distinguish receipt of evidence from execution when the latter is unknown.
This version includes evidence available through \evidencedate.
